% Generated from the matching Typst scaffold by
% scripts/migrate_typst_report_to_latex.py.

% ── 2.2  The Plasticity–Stability Tradeoff and Forgetting ────────────────────
% PURPOSE: Establish the stability/plasticity tension and the counter-intuitive thesis
%          that forgetting can HELP, motivating proactive memory policies.
% MOVE: Counter-intuitive lead — contrary to "more memory is always better", accumulated
%       irrelevant experience can degrade an agent.
% BEATS: see bullets below.
% FIGURES/TABLES: (none — narrative; forward-ref fig:interdependence for error propagation)
% CITATIONS: xiong2025 (experience-following, error propagation), [VERIFY+ADD] Context-Rot (Chroma)
% ─────────────────────────────────────────────────────────────────────────────

\section{The Plasticity--Stability Tradeoff and Forgetting}

\begin{todobox}
Expand the beats below into prose using the Counter-intuitive lead move; land on the claim that proactive forgetting is a design lever, not merely a failure to retain.
\end{todobox}

% - Stability vs plasticity: retaining prior competence pulls against adapting to new tasks;
%   an unbounded store maximizes stability but risks drowning the relevant in the irrelevant.
% - Catastrophic forgetting is the classic failure mode — but it is framed at the weight level;
%   for external memory the analogous risk is accumulation, not erasure.
% - Counter-intuitive case: forgetting can HELP — old or off-task memory acts as a distraction
%   that crowds out the relevant item and biases retrieval (recurring theme: "footprint").
% - Error propagation: agents exhibit experience-following — a bad stored experience is
%   retrieved, reused, and compounds across subsequent tasks @xiong2025.
% - Hence proactive forgetting is a control lever; thread "not every task needs memory" and
%   forward-ref the interdependence analysis (fig:interdependence) where ~⅔ of tasks are
%   file-disjoint. [VERIFY+ADD: Context-Rot (Chroma) — long irrelevant context degrades quality]
